\section{Related Work}
\label{sec:related}

We review three strands of prior work used below: DMDGP partial-reflection symmetry, DDGP recognition, and the negative counting result for general DDGP instances.

In the DMDGP, the predecessor sets $U_j$ are contiguous and immediately precede $j$ in the vertex order, so $U_j = \{j-K, \ldots, j-1\}$~\cite{lavor2012discretizable}. This contiguity ensures that each $U_j$ induces a clique of size $K$ in $G$. The clique property implies that the branch choices at each level are partial reflections across affine hyperplanes spanned by predecessor coordinates~\cite{mucherino2012exploiting}. These reflections form a group of partial-reflection symmetries acting on the BP tree. They explain the generic power-of-two solution counts and the structure of the search space~\cite{liberti2011number,liberti2014number}. 

Because these reflections act on suffix intervals of the vertex order, one can count the number of feasible BP paths by a combinatorial method depending on graph topology rather than numerical edge weights~\cite{lavor2021optimality}. This counting method has been used to prove fixed-parameter tractability of the BP algorithm on DMDGP instances with respect to $K$ and a branching-frequency parameter $v_0$. Partial-reflection symmetries also support build-up, splitting, and search algorithms for DMDGP instances~\cite{mucherino2012exploiting,goncalves2021new}. 

Finding a valid DMDGP vertex order from a graph is NP-complete for any fixed dimension $K \ge 2$~\cite{cassioli2015discretization}. In the general DDGP, the predecessor sets $U_j$ need not be consecutive. This relaxation changes the computational tradeoff: finding a valid DDGP order can be done in polynomial time by an FPT algorithm with respect to $K$~\cite{abud2024impossible}. 

The same relaxation complicates counting. Since DDGP predecessors are not necessarily consecutive, reflections no longer act on suffix intervals. A branch change at a vertex $q$ propagates through a directed dependency graph of successors. This dynamic propagation is the obstruction we address in the present local construction.

Abud et~al.~\cite{abud2024impossible} proved that no purely combinatorial counting method exists for the general DDGP. Their 5-vertex counterexample has a predecessor set $U_5 = \{1,4\}$ that does not induce a clique. After branching on vertex 4, one candidate coordinate for vertex 5 can violate the triangle inequality on $\{1,4,5\}$ with positive probability. Thus the number of valid coordinates for vertex 5, and therefore the total solution count, depends on the numerical edge weights rather than only on graph topology. Our result does not contradict this barrier: it gives a generic rank count for a fixed local active subproblem, not a global combinatorial count for arbitrary DDGP instances.
